\documentclass[12pt]{article}
\usepackage{graphicx}
\usepackage{amsmath}
\usepackage{booktabs}
\usepackage{enumitem}
\usepackage[a4paper,margin=1in]{geometry}
\usepackage{array}

\begin{document}


\section*{Question 1}

Enzymes are catalysts. \\
(a) (i) Describe what is meant by the term catalyst. \hfill [2] \\
\noindent\rule{15cm}{0.4pt} \\[3em]

(ii) Enzymes are proteins. State the names of the chemical elements found in all proteins. \hfill [2] \\
\noindent\rule{15cm}{0.4pt} \\[3em]

\section*{Question 2}

Lipase is an enzyme found in the human digestive system. \\
Fig.~1 shows the effect of pH on the percentage activity of lipase:

\begin{center}
\includegraphics[width=0.8\textwidth]{images/Fig1_1_LipaseActivity.png}
\end{center}
\textbf{Fig.~1.} Effect of pH on lipase activity.

Describe and explain the change in percentage activity of lipase shown on Fig.~1. \hfill [6] \\
\noindent\rule{15cm}{0.4pt} \\[4em]

\section*{Question 3}

Fig.~2 shows a cross-section of a human heart.

\begin{center}
\includegraphics[width=0.55\textwidth]{images/Fig2_1_HeartCrossSection.png}
\end{center}
\textbf{Fig.~2.} Cross-section of the human heart.

(a) On Fig.~2, label with an R the position of the right ventricle. \hfill [1] \\
\noindent\rule{10cm}{0.4pt} \\[2em]

(ii) State the names of the structures labelled A and B on Fig.~2. \hfill [2] \\
A: \rule{8cm}{0.4pt} \\
B: \rule{8cm}{0.4pt} \\[1em]

(iii) Complete the sentence about blood vessel C. Blood vessel C transports blood from the \underline{\hspace{5cm}} to the \underline{\hspace{5cm}}. \hfill [2] \\[2em]

\section*{Question 4}

Table 5.1: Plants growing in Europe

\begin{center}
\begin{tabular}{cc}
\toprule
Plant Number & Leaf Area (cm$^2$) \\
\midrule
1 & 310 \\
2 & 335 \\
3 & 390 \\
4 & 348 \\
5 & 365 \\
\midrule
Mean & 350 \\
\bottomrule
\end{tabular}
\end{center}

(i) Using the data in Table 5.2, calculate the mean leaf area for plants growing in North America. Give your answer as a whole number and include the unit. \hfill [2] \\
\noindent\rule{12cm}{0.4pt} \\[3em]

(ii) State a conclusion for the results shown in Table  ............................................ \hfill [1] \\

(iii) Ancestors of monkey flower plants had a much smaller leaf area. \\
Explain how the monkey flower plants have developed a larger leaf area over time. \hfill [5] \\
\noindent\rule{15cm}{0.4pt} \\[3em]

State two advantages of using a pyramid of energy rather than a pyramid of biomass to represent a food chain. \hfill [2] \\

\section*{Question 5}

(a) The oceans contain fish stocks that can be managed as a sustainable resource to provide food for humans. \\

(i) State what is meant by the term sustainable resource. \hfill [2] \\

(ii) Fishing can be managed to promote the conservation of fish stocks. \\
Explain one way that fishing can be managed to conserve fish stocks. \hfill [2] \\



% Question 1
\noindent\textbf{1 Biotechnology is used in the process of bread-making.}

\begin{center}
\includegraphics[width=0.7\textwidth]{images/bread_making_steps.png}
\end{center}

\noindent (a) State the name of organism A in the figure above. [1] \\
\vspace{1.5cm}

\noindent (b) (i) State the name of the source of energy used by organism A. [1] \\
\vspace{1.5cm}

\noindent (ii) State the name of the process occurring at step 3 that causes gas bubbles to form in the dough. [1] \\
\vspace{1.5cm}

\noindent (iii) State the name of the gas that forms to create the gas bubbles in the dough. [1] \\
\vspace{1.5cm}

\noindent (c) Explain the reasons for the different temperatures used in step 3 and step 5. [2] \\
\vspace{3cm}

\noindent (d) State the names of two products of biotechnology, other than bread, that make use of microorganisms. [2] \\
\vspace{3cm}

\newpage

% Question 2
\noindent\textbf{2 The concentration of atmospheric carbon dioxide has increased considerably in recent years.}

\noindent (a) Describe the possible causes of increased atmospheric carbon dioxide. [3] \\
\vspace{4cm}

\noindent (b) Soybean plants, Glycine max, were grown in two plots with different CO$_2$ concentrations, 370 ppm and 550 ppm. Fig. 2.1 shows average photosynthesis rates.

\begin{center}
\includegraphics[width=0.8\textwidth]{images/photosynthesis_rate_graph.png}
\end{center}

\noindent Describe and explain the effect of carbon dioxide concentration on the average rates of photosynthesis of the soybean plants from 04:00 to 22:00. [6] \\
\vspace{5cm}

\noindent (c) Complete Table 2.1 by stating two structural features of epidermis and mesophyll tissues and explain how each feature is an adaptation for photosynthesis. [4] \\
\vspace{4cm}

\noindent (d) Scientists had higher breathing rates in the 550 ppm CO$_2$ plot. Suggest why. [2] \\
\vspace{2cm}

\newpage

% Question 3
\noindent\textbf{3 Very small pieces of plastic called microplastics are found in many products such as soaps and toothpaste.}

\noindent (a) (i) It is estimated that microplastics make up 5\% of some toothpastes by mass. Each person uses 2 g of toothpaste daily. In 2013, 1.2$\times10^{9}$ people used toothpaste containing microplastics.

Calculate the mass of microplastics used on one day in 2013 in these countries. Show working and units. [3] \\
\vspace{3cm}

\noindent (ii) State one recommendation, other than regular brushing, for proper care of teeth. [1] \\
\vspace{1.5cm}

\noindent (b) Lugworms live in sand on beaches and feed on diatoms, which require ammonium ions. Researchers tested effects of microplastics on lugworms; see Table 3.1.

Describe and explain why researchers concluded that non-biodegradable microplastics are most harmful to lugworms using Table 3.1. [4] \\
\vspace{5cm}

\noindent (c) State the names of molecules converted into ammonium ions in the nitrogen cycle. [1] \\
\vspace{1.5cm}

\noindent (d) State the name of the process converting ammonium ions into nitrate ions. [1] \\
\vspace{1.5cm}

\noindent (e) Explain the effects of nitrate ion deficiency on plant growth. [3] \\
\vspace{3cm}

\noindent (f) Discuss the effects of non-biodegradable plastics on terrestrial ecosystems. [5] \\
\vspace{4cm}

\newpage

% Question 4
\noindent\textbf{4 Neurones are part of the nervous system connected by synapses.}

\noindent (a) Describe how the structure of a neurone relates to its function. [3] \\
\vspace{3cm}

\noindent (b) State the organs that make up the central nervous system. [1] \\
\vspace{1.5cm}

\noindent (c) Outline the reflex arc pathway in response to shining bright light into the eye. [3] \\
\vspace{3cm}

\noindent (d) Suggest why reflexes occur in unconscious people. [1] \\
\vspace{1.5cm}

\noindent (e) shows a synapse and parts of two neurones.

\begin{center}
\includegraphics[width=0.6\textwidth]{images/synapse_diagram.png}
\end{center}

(i) Draw an arrow to show the direction of the signal across the synapse. [1] \\
\vspace{3cm}

\newpage

% Question 5
\noindent\textbf{5 The testes are part of the endocrine and reproductive systems.}

\noindent (a) State the hormone released from the testes. [1] \\
\vspace{1.5cm}

\noindent (b) Complete Fig. 5.1 by stating two other organs that belong to two organ systems.

\begin{center}
\includegraphics[width=0.5\textwidth]{images/organs_two_systems.png}
\end{center}
\vspace{1.5cm}

\noindent (c) Fig. 5.2 shows photomicrograph of part of mammalian testis.

\begin{center}
\includegraphics[width=0.7\textwidth]{images/mammalian_testis_micrograph.png}
\end{center}

Explain why meiosis is necessary in the testes. [3] \\
\vspace{3cm}

\noindent (d)  shows a photomicrograph of a sperm cell.

\begin{center}
\includegraphics[width=0.7\textwidth]{images/sperm_cell_micrograph.png}
\end{center}

\vspace{3cm}

\noindent (e) Draw and label one human egg cell, including at least one feature not found in a sperm cell. [3] \\
\vspace{3cm}

\noindent (f) Describe what happens to a fertilised egg cell before implantation in the uterus. [3] \\
\vspace{3cm}

\newpage

% Question 6
\noindent\textbf{6 Fig. 6.1 shows some varieties of potato cultivated worldwide.}

\begin{center}
\includegraphics[width=0.6\textwidth]{images/potato_varieties.png}
\end{center}

\noindent (a) Define the term ‘species’. [2] \\
\vspace{2cm}

\noindent (b) Fig. 6.2 shows a reproduction method used by some potato farmers.

\begin{center}
\includegraphics[width=0.6\textwidth]{images/potato_tubers_method.png}
\end{center}

Describe the advantages of this type of reproduction in crop production. [3] \\
\vspace{3cm}

\noindent (c) Explain why plants store starch in potato tubers. [2] \\
\vspace{2cm}

\end{document}
